\documentclass{article}
\usepackage{amsmath}
\usepackage{amssymb}
\usepackage{amsthm}


\newtheorem*{conjecture}{Conjecture}
\newtheorem*{theorem}{Theorem}
\newtheorem*{lemma}{Lemma}


\makeatletter
\newcommand*{\rom}[1]{\expandafter\@slowromancap\romannumeral #1@}
\makeatother


\begin{document}

\section*{Exercise 1.2.1}

\subsection*{(a)}

True, see property VS 3

\subsection*{(b)}

False.

Suppose \(0_{1}\) and \(0_{2}\) are both zero vectors in a vector space \(V\) over a field \(F\). It follows from (VS 3) and (VS 1):

\begin{equation*}
    0_{1} = 0_{1} + 0_{2} = 0_{2} + 0_{1} = 0_{2}
\end{equation*}

And \(0_{1} = 0_{2}\) implies that a vector space cannot have more than one zero vector
(and it is unique).

\subsection*{(c)}

False

Assume a vector space \(V\) over a field \(F\). Suppose \(x = 0 \in V\) and
\(a \neq b\) both in \(F\). We have now:

\begin{equation*}
    ax = 0 = bx 
\end{equation*}

but \(a \neq b\).

\subsection*{(d)}

False

Assume a vector space \(V\) over a field \(F\). Suppose \(x \neq y\) are vectors
in \(V\). Let \(a = 0 \in F\). We have:

\begin{equation*}
    ax = 0 = ay
\end{equation*}

but \(x \neq y\)

\subsection*{(e)}

True, Axler is way more fun :(

\subsection*{(f)}

False, see the definition in Example 1 (Down with determinants!)

\subsection*{(g)}

False (I push my fingers into my ... EYESSSSSSSSSSSSSSSSSSS)

Suppose \(m < n\). Let

\begin{equation*}
    f(x) = a_{n}x^{n} + a_{n-1}x^{n-1} + \dots + a_{1}x + a_{0}
\end{equation*}
\begin{equation*}
    g(x) = b_{m}x^{m} + b_{m-1}x^{m-1} + \dots + b_{1}x + b_{0}
\end{equation*}

We can always write \(g(x)\) using the index \(n\) equivalently, by making the terms above \(m\) zero, writing:

\begin{equation*}
    g(x) = 0x^{n} + 0x^{m+2} + 0x^{m+1} + b_{m}x^{m} + \dots + b_{1}x + b_{0}
\end{equation*}

Which can be written more clearly:

\begin{equation*}
    g(x) = b_{n}x^{n} + b_{n-1}x^{n-1} + \dots + b_{1}x + b_{0}
\end{equation*}

by setting:

\begin{equation*}
    0 = b_{n} = \dots = b_{m+2} = b_{m + 1}
\end{equation*}

\subsection*{(h)}

True, why not Axler micro :(

\subsection*{(i)}

True, why not Axler aleps?

\subsection*{(j)}

True, I love you both, we're totally doing Friedberg et. al!

\subsection*{(k)}

True

\section*{Exercise 1.2.8}

Suppose \(V\) is a vector space over a field \(F\). Assume \(x, y \in V\) and \(a, b \in F\). From 
the properties of linear spaces:

\begin{equation*}
    (a + b)(x + y) = (a + b)x + (a + b)y \ \ \ \emph{Using  VS 7}
\end{equation*}
\begin{equation*}
    = ax + bx + ay + by \ \ \ \emph{Using VS 8}
\end{equation*}
\begin{equation*}
    = ax + ay + bx + by \ \ \ \emph{Using VS 1}
\end{equation*}

and this completes the proof.

\end{document}